\paragraph{应用。}\label{sub:application}
我们遵循与高维情形相同的策略。取\trDoubleBranchedCover{}$\pi\colon X=X^1_m\to\CC$ 的\trVanishingPath{}的一组基，并考虑 $\scr{F}(\pi)$ 中与之相伴的\trFullSubcategory，我们把它记作 $\scr{B}^{\rightarrow}$。明确地说，它具有\trObject{}$(Z_0,\ldots,Z_m)$，使得
\begin{equation*}
\hom_{\scr{B}^{\rightarrow}}(Z_i,Z_j)=
\begin{cases}
\KK\cdot c_{ji}\oplus\KK\cdot d_{ji}, & i<j,\\
\KK\cdot e_{ii}, & i=j,\\
0, & i>j,
\end{cases}
\end{equation*}
其中所有\trDegree{}均为零。诸 $e_{ii}$ 是严格的单位元，除此之外我们还有
\begin{equation*}
\begin{gathered}
\mu^2_{\scr{B}^{\rightarrow}}(c_{kj},c_{ji})=c_{ki},\qquad
\mu^2_{\scr{B}^{\rightarrow}}(d_{kj},d_{ji})=d_{ki},\\
\mu^2_{\scr{B}^{\rightarrow}}(d_{kj},c_{ji})=0,\qquad
\mu^2_{\scr{B}^{\rightarrow}}(c_{kj},d_{ji})=0.
\end{gathered}
\end{equation*}
这里与高维类比 \eqref{eq:directed-category-hom} 之间的差异，最终可以归结为这样一个事实：$S^0$ 的\trCohomology{}环与高维\trSphere{}的表现不同。取由
\begin{equation}\label{eq:twisted-complexes-n1}
C_0=Z_0,\qquad C_j=\mathrm{Cone}(c_{j,j-1}\pm d_{j,j-1})\quad\text{（当 } j>0\text{ 时）}
\end{equation}
给出的\trTwistedComplex{}$(C_0,\ldots,C_m)$——其中每个 $j$ 处的符号都可以自由选取（关于这种歧义性的讨论见下面的\cref{thm:sign-ambiguity-remark}）——并设 $\scr{A}$ 是 $Tw\,\scr{B}^{\rightarrow}$ 中与之相伴的\trFullSubcategory。容易算出 $H(\scr{A})\cong\bar{A}$。
\begin{lemma}\label{thm:extended-not-formal}
$\bar{\scr{A}}$ 不是形式的。
\end{lemma}
\begin{proof}
显然只需对 $m=2$ 证明这一点，因为其余的 $\bar{\scr{A}}$ 都把这一个作为\trFullSubcategory{}包含在内。由于族 \eqref{eq:twisted-complexes-n1} 生成 $Tw(\scr{B}^{\rightarrow})$，$\scr{A}$ 与 $\scr{B}^{\rightarrow}$ 上的\trTwistedComplex{}的\trACategory{}是\trQuasiequivalent{}的，而这又蕴含它们具有\trIsomorphic{}的\trHochschildCohomology{}群（这类结果可追溯到 [71] 对代数的处理，其证明稍作调整便可在\trAinfty{}的情形中通过）。直截了当地计算给出
\begin{equation*}
HH(\scr{B}^{\rightarrow},\scr{B}^{\rightarrow})=\KK\ \text{（位于零次）}\oplus\KK^2\ \text{（位于一次）}
\end{equation*}
这与\cref{thm:extended-hh-pieces}中（部分地）算出的 $\bar{A}$ 的\trHochschildCohomology{}不符，所欲的非形式性结果便由此得证。
\end{proof}
设 $\scr{A}$ 是 $\bar{\scr{A}}$ 中以 $(C_1,\ldots,C_m)$ 为\trObject{}的\trFullSubcategory。作为前面讨论的一个后果，这又将是非形式的，至少对 $m\ge 5$ 如此。另一方面，\cref{thm:matching-l-cone}（经适当扩充，见\cref{thm:sign-ambiguity-remark}）表明 $\scr{A}$ \trQuasiisomorphic{}于 $\scr{F}$。于是，
\begin{corollary}\label{thm:f1-not-formal}
$\scr{F}=\scr{F}^1_m$ 对任何 $m\ge 5$ 都不是形式的。
\noqed
\end{corollary}
\begin{remark}\label{thm:sign-ambiguity-remark}
我们曾在二维情形中使用\cref{thm:matching-l-cone}，在那里\trLefschetzFibration{}约化成一个\trDoubleBranchedCover。尽管这种情形未包含在我们原来的讨论之中，其证明却不难适应，或者也可以改用一个更初等的论证来代替。然而，在陈述精确的命题时需要一些小心（这一点在\cref{thm:zero-dimensional-extension}中已简要提到过）。\trMorphism{}$c_{j,j-1}$ 与 $d_{j,j-1}$ 对应于诸\trLefschetzThimble{}之间的交点，而结论是：与之相伴的\trMatchingCycle{}是 $c_{j,j-1}\pm d_{j,j-1}$ 上的\trCone，符号的选取决定了该\trMatchingCycle{}携带哪个\trPinStructure。在我们的情形，诸\trMatchingCycle{}是 $X$ 上构成 $(A_m)$ 构形的那些曲线，而我们原来假定它们携带非平凡的\trPinStructure（因为我们曾应用\cref{thm:twist-equals-tau}证明了它们\trSplitGeneration）。幸运的是，这些曲线恰好在 $H_1(X;\ZZ/2)$ 中\trLinear{}无关，这意味着任何两种\trPinStructure{}的选取都可以通过用 $X$ 上的一个\trRealLineBundle{}作扭转而相互得到（用\cref{sub:automorphisms}的语言说，就是一个\trBrane\trAutomorphism）。代数上，这反映在 $\scr{B}^{\rightarrow}$ 容许改变 $c_{j,j-1}\pm d_{j,j-1}$ 之符号的\trAutomorphism{}这一事实之中。这正是 \eqref{eq:twisted-complexes-n1} 中符号的选取无关紧要的原因。
\end{remark}
